\chapter{系统需求分析}

\newcommand{\chapterthreeplaceholder}[1]{%
\begin{samepage}
\begin{center}
\fbox{\parbox[c][4.0cm][c]{0.80\textwidth}{\centering #1（此处插入）}}
\end{center}
\begin{center}
#1（此处插入）
\end{center}
\end{samepage}
}

\newcounter{usecaseitem}[subsection]
\newcommand{\usecasesubheading}[1]{%
  \refstepcounter{usecaseitem}%
  \vspace{1ex}
  \noindent\hspace*{2em}{\MiyaSubsubSecfont (\arabic{usecaseitem})\ #1}\par
}
\newcommand{\circlednum}[1]{%
  \ifcase#1\or \ding{192}\or \ding{193}\or \ding{194}\or \ding{195}\or \ding{196}\or \ding{197}\or \ding{198}\or \ding{199}\or \ding{200}\or \ding{201}\else #1\fi
}
\newlist{usecaselist}{enumerate}{1}
\setlist[usecaselist]{label=\protect\circlednum{\arabic*}, leftmargin=2.8em, labelsep=0.6em, itemsep=0.2ex, topsep=0.4ex}

\section{需求概述}

本章围绕 Mini-Chat-Bar 面向的实际使用场景展开需求分析。本文关注的不是一般社交聊天，而是技术讨论中的高频操作，例如登录后进入私聊或聊天室、发送代码和文件、回看历史消息、对关键内容做收藏，以及在无人及时回复时调用 AI 继续处理问题。第 1 章已经说明，这类场景对上下文连续性、消息组织方式和讨论结果沉淀都有更高要求，因此需要把需求拆得更细，不能只停留在实现聊天功能这一层。

结合当前系统定位，需求分析主要围绕用户对象、使用场景、核心功能以及性能、安全和可维护性要求展开。后文将按照业务需求、功能需求、非功能需求和用户用例四个部分展开。

\section{业务需求分析}

业务需求分析的目标是先把“谁在用、怎么用、为什么用”说清楚。只有用户、场景和流程明确，后续的模块设计和实现细节才有依据。

\subsection{用户分析}

根据本系统的使用目标，用户可以概括为技术交流人群、个人开发者和普通交流用户三类。三类用户都需要基础聊天能力，但关注重点并不相同。

\subsubsection{技术交流人群}

技术交流人群包括课程项目组成员、实验室小团队、开源协作成员等。这类用户最常见的操作是在会话中发送代码片段、报错日志、接口参数和测试结果，然后围绕同一个问题持续讨论。他们关注的不是普通聊天，而是问题能不能尽快定位、结论能不能保存、后来的人能不能看懂前面的讨论。因此，这类用户对聊天室、历史回读、引用回复、文件共享和 AI 辅助问答的依赖最强，这类面向问题定位的持续讨论方式，与面向任务协作的即时交流系统研究结论较为一致\cite{imdesign2023}，使用流程如图~\ref{fig:3-1} 所示。个人开发者和普通交流用户的流程与该主链路较为接近，因此在后文通过典型场景和用例分别说明，不再重复绘制结构相近的流程图。

\begin{figure}[H]
\centering
\thesisfigwide{img/图3_1.png}
\caption{技术交流人群使用流程示意图}
\label{fig:3-1}
\end{figure}

\subsubsection{个人开发者}

个人开发者通常独立完成学习或开发任务，遇到问题时更倾向于先自己查资料、再向系统提问。他们对侧边栏 AI 助手、历史检索和收藏功能的需求更明显：一方面，希望快速得到代码解释、报错分析或实现建议；另一方面，希望把有价值的回答和讨论沉淀下来，形成个人资料库。当单人问答不足以解决问题时，这类用户也会临时创建聊天室，与其他开发者集中讨论。

\subsubsection{普通交流用户}

普通交流用户可以是测试人员、产品人员或只需要跟进阶段进展的参与者。这类用户并不一定频繁发起技术讨论，但需要较低门槛地进入房间、查看阶段性结论、下载共享文件和回读历史消息。他们更看重界面是否直观、总结是否清楚，以及能否在不参与完整讨论过程的情况下快速了解当前进展。

本节以技术交流人群作为核心目标用户展示流程图。这类用户的操作链路最完整，能够把聊天室协同、AI 补充分析、结果收藏与后续复用等关键环节连成一条主线，更能体现系统围绕“技术讨论”展开的设计重点。
\subsection{业务场景描述}

结合上述用户特征，可以归纳出本系统中较为典型的三类业务场景。

\subsubsection{场景一：多人协同定位接口异常}

项目成员在聊天室中集中处理接口异常问题。发起者先发送报错信息、请求参数和相关代码，其他成员在同一房间继续补充判断；如果短时间内没有形成明确结论，可以直接在房间中触发 \texttt{@AI} 进行辅助分析。问题解决后，关键回答和结论还应允许被收藏，方便后续再次查阅。

\subsubsection{场景二：个人开发者快速验证实现思路}

个人开发者在实现某个功能时，可能先通过 AI 面板询问框架用法、错误原因或代码修改建议。如果单轮问答不能解决问题，用户可以新建一个与当前主题相关的聊天室，把问题发给其他成员继续讨论。讨论结束后，重要结论可以保存在收藏页中，避免下次遇到同类问题时重复检索。

\subsubsection{场景三：普通交流用户整理阶段性讨论结论}

测试人员或产品人员未必全程参与讨论，但在查看功能变更、确认交付内容或准备汇报材料时，需要快速掌握前面的讨论结果。此时他们更依赖聊天总结、历史回读和共享文件下载等功能，以减少重新沟通的成本。

\subsection{业务流程分析}

从系统角度看，一次完整的业务流程通常包括登录进入、选择会话、发送或接收消息、触发 AI 辅助、整理关键信息和后续回看几个阶段，如图~\ref{fig:3-2} 所示。

\begin{figure}[H]
\centering
\thesisfigflow{img/图3_2.png}
\caption{系统业务流程图}
\label{fig:3-2}
\end{figure}

用户首先完成注册或登录，进入聊天主界面后选择私聊或聊天室。在交流过程中，系统需要支持文本、代码、图片和文件等多种消息形式；如果讨论涉及持续协作，还需要显示房间成员和在线状态。遇到技术问题时，用户可以在侧边栏 AI 面板单独提问，也可以在聊天室中直接发起 \texttt{@AI} 提问。讨论结束后，重要消息可以加入收藏，聊天室消息和 AI 回复应允许通过历史记录再次读取；对于设置了有效期的聊天室，任务结束后还应转为归档状态，供后续查阅。

\section{功能需求分析}

根据业务需求，系统功能可归纳为用户管理、即时通讯、聊天室、AI 助手和收藏五个核心模块，具体功能汇总见表~\ref{tab:3-1}。各模块之间的结构关系将在第4章系统设计部分结合功能结构图说明。

\begin{table}[!htbp]
\centering
\caption{系统功能需求汇总}
\label{tab:3-1}
\small
\renewcommand{\arraystretch}{1.25}
\setlength{\tabcolsep}{1ex}
\resizebox{\textwidth}{!}{%
\begin{tabular}{p{0.14\textwidth}p{0.22\textwidth}p{0.38\textwidth}p{0.16\textwidth}}
\toprule
\textbf{功能模块} & \textbf{核心功能点} & \textbf{功能说明} & \textbf{主要服务对象} \\
\midrule
用户管理模块 & 注册登录、资料维护、在线状态 & 完成身份认证、用户资料读取与更新，并为后续会话和权限控制提供基础数据 & 全部用户 \\
即时通讯模块 & 私聊、消息发送、历史回读、富媒体支持 & 承担私聊消息收发与历史回看，支持文本、代码、图片和文件等消息类型 & 全部用户 \\
聊天室模块 & 房间创建、加入方式、成员状态、归档管理 & 提供面向技术讨论的多人协作空间，支持公开或密码加入、在线人数同步和到期归档 & 技术交流人群、个人开发者 \\
AI 助手模块 & 单人问答、房间问答、总结生成、上下文检索 & 在侧边栏或聊天室内提供技术问答、代码解释和讨论总结能力，并结合历史消息提升回答相关性 & 技术交流人群、个人开发者 \\
收藏模块 & 收藏、分类、检索、回看来源 & 将关键消息、代码或结论保存下来，支持后续检索与再次回看 & 全部用户 \\
\bottomrule
\end{tabular}
\par}
\end{table}

\subsection{用户管理模块}

用户管理模块负责系统的身份入口和资料基础，既要处理注册与登录，也要为会话页、聊天室和收藏页提供稳定的用户身份数据。

\begin{enumerate}[label={(\arabic*)}, leftmargin=2.4em]
\item 支持邮箱密码登录、邮箱验证码登录以及第三方登录，保证不同用户可以顺利进入系统。
\item 登录成功后生成有效会话信息，并在页面刷新或重新进入时保持登录状态。
\item 支持读取和修改昵称、头像等基础资料，使聊天列表、消息记录和个人中心中的信息保持一致。
\item 支持在线状态同步，为联系人列表和聊天室成员列表提供状态依据。
\end{enumerate}

\subsection{即时通讯模块}

即时通讯模块是系统的核心链路，主要负责私聊消息的发送、接收、展示和历史回读。由于本系统面向技术交流，消息类型不能只限于普通文本。

\begin{enumerate}[label={(\arabic*)}, leftmargin=2.4em]
\item 支持用户之间进行私聊，并能够及时同步新消息。
\item 支持文本、代码片段、图片、文件等多种消息形式，保证技术讨论内容能够被完整表达。
\item 支持引用回复、消息状态反馈和历史消息分页读取，便于围绕具体内容继续讨论。
\item 支持最近会话、未读消息和消息搜索等基础能力，降低回看和定位信息的成本。
\end{enumerate}

\subsection{聊天室模块}

聊天室模块主要服务于多人协作讨论，适合处理某个专题、某次发布或某个待解决问题。相比私聊，聊天室更强调成员组织、消息广播和阶段性归档。

\begin{enumerate}[label={(\arabic*)}, leftmargin=2.4em]
\item 支持创建聊天室，并设置名称、技术方向、加入方式和有效期等基础信息。
\item 支持公开加入和密码加入等方式，满足不同讨论范围的需要。
\item 支持成员进入房间后的在线人数同步、历史消息回读和新消息广播。
\item 支持聊天室到期后转为归档状态，避免长期闲置房间继续产生无效消息。
\end{enumerate}

\subsection{AI 助手智能问答模块}

AI 助手模块负责补充人工讨论中的空档期，帮助用户更快获得初步答案。该模块既要支持单人连续追问，也要支持结合聊天记录的多人上下文问答，并承担术语解释与阶段总结等辅助任务，这与对话历史管理型问答系统的设计重点相吻合\cite{historyragqa2026}。

\begin{enumerate}[label={(\arabic*)}, leftmargin=2.4em]
\item 支持侧边栏 AI 助手，供用户单独输入问题、代码或报错信息，并对选中的术语、文本或代码片段进行解释。
\item 支持在聊天室中通过 \texttt{@AI} 触发问答，使房间成员可以共同查看 AI 返回结果。
\item 支持结合历史消息进行检索增强，使 AI 回答尽量贴合当前讨论上下文，减少重复说明前情的成本。
\item 支持根据一段时间内的聊天内容生成讨论总结，便于后来加入的成员快速了解前情。
\end{enumerate}

\subsection{收藏模块}

收藏模块用于把聊天过程中的临时信息转化为可回看的长期资料，这也是本系统区别于普通聊天工具的重要部分。

\begin{enumerate}[label={(\arabic*)}, leftmargin=2.4em]
\item 支持对私聊消息、聊天室消息和 AI 回复执行收藏操作。
\item 支持为收藏内容增加备注或标签，方便后续按主题整理。
\item 支持在收藏页中查看已保存内容，并保留来源信息，便于回到原始会话继续阅读。
\item 支持基础检索与筛选能力，帮助用户快速定位历史问题和解决方案。
\end{enumerate}

\section{非功能需求分析}

除业务功能外，系统还需要满足基本的性能、安全、可用性和可维护性要求。这些要求虽然不直接对应某个按钮或页面，但会直接影响系统能否稳定运行。

\subsection{性能需求}

本系统主要面向课程项目、小型技术团队和答辩演示场景，因此性能要求以稳定可用为主，而不是追求大规模平台级指标。普通业务接口如登录、资料查询、消息列表读取和收藏列表读取应保持较快响应，常见查询在教学演示规模下应控制在秒级以内；私聊和聊天室消息发送后需要及时回显，在线用户应在可感知实时范围内收到更新。对于 AI 功能，由于依赖外部模型服务，其耗时可以明显高于普通接口，但界面需要提供思考状态或流式反馈，避免用户误以为系统失去响应\cite{restfulapitest2025}。

\subsection{安全性需求}

系统需要保证身份认证、访问控制和数据传输的基本安全。用户密码应采用单向散列方式保存，避免明文存储；登录后所有受保护接口都应进行身份校验，并限制未授权用户访问他人资料、私聊记录和受限聊天室。对于聊天室加入、文件上传和 AI 上下文读取，也需要结合当前用户权限做边界检查，防止越权读取或恶意输入影响系统稳定性\cite{imsecurity2024}。

\subsection{可用性需求}

系统界面应保持清晰直接，避免把高频操作埋得过深。登录、切换会话、发送消息、查看历史、触发 AI 和收藏内容都应在较少步骤内完成。对于技术交流场景来说，代码、文件和普通文本消息需要有明显区分，避免阅读混乱；对于桌面端和 Web 端，还应保证基础操作逻辑尽量一致，减少用户在不同端之间切换时的学习成本。

\subsection{数据分析需求}

虽然本系统不是数据分析平台，但为了后续优化系统运行和用户体验，仍需要保留一些基础统计项。这些数据主要服务于后续观察系统负载、房间使用情况、AI 调用频率和收藏使用情况，不直接面向普通用户展示，核心数据项见表~\ref{tab:3-2}。

\begin{table}[!htbp]
\centering
\caption{系统数据分析核心数据项}
\label{tab:3-2}
\small
\renewcommand{\arraystretch}{1.25}
\setlength{\tabcolsep}{1ex}
\resizebox{\textwidth}{!}{%
\begin{tabular}{p{0.17\textwidth}p{0.24\textwidth}p{0.33\textwidth}p{0.16\textwidth}}
\toprule
\textbf{数据类别} & \textbf{关键数据项} & \textbf{采集目的} & \textbf{更新方式} \\
\midrule
系统活跃数据 & 在线用户数、连接持续时间 & 观察系统负载和在线波动情况 & 实时记录/按日汇总 \\
聊天室使用数据 & 房间创建数、过期房间数、成员活跃度 & 判断聊天室功能的实际使用情况 & 实时记录/按周汇总 \\
消息流量数据 & 私聊消息量、聊天室消息量、文件消息量 & 估算消息存储和网络传输压力 & 实时记录/按日汇总 \\
AI 调用数据 & 提问次数、总结次数、响应耗时 & 评估 AI 模块使用频率和调用开销 & 调用时记录/按日汇总 \\
收藏沉淀数据 & 收藏总量、标签使用情况、检索频率 & 判断知识沉淀功能的利用效果 & 实时记录/按周汇总 \\
\bottomrule
\end{tabular}
\par}
\end{table}

\subsection{可扩展性与可维护性需求}

系统应保持清晰的模块边界，便于后续增加新消息类型、替换 AI 服务或扩展桌面端功能。前后端分离结构、接口分层和独立的控制器/路由设计有助于降低修改时的耦合度；同时，系统还应保留必要的日志与错误信息，便于联调、测试和后续维护。

\section{用户用例分析}

在需求分析阶段，用例的作用是从用户视角检验系统是否覆盖了主要操作路径。结合前文的用户分类，下面分别说明三类用户的典型用例。前文的图~\ref{fig:3-1} 已经用完整流程图突出技术交流人群这一核心目标用户，图~\ref{fig:3-2} 则概括了三类用户共享的系统业务闭环。在此基础上，本节针对个人开发者和普通交流用户，分别给出代表性用例，重点分析他们在 AI 求解、历史回看、文件获取和结果沉淀四个环节中的使用行为

\subsection{技术交流人群用例分析}

技术交流人群的使用路径最完整，通常会同时涉及聊天室、消息回读、AI 提问和收藏等多个功能，用例如图~\ref{fig:3-3} 所示。

\begin{figure}[H]
\centering
\thesisfigflow{img/图3_4.png}
\caption{技术交流人群用例图}
\label{fig:3-3}
\end{figure}

\usecasesubheading{用例一：发起技术问题求助}

当项目成员在开发或联调过程中遇到问题时，可按以下流程完成求助：
\begin{usecaselist}
\item 进入对应聊天室，发送问题描述、报错日志和相关代码片段。
\item 其他成员在房间内继续补充分析，必要时通过 \texttt{@AI} 获取辅助回答。
\item 问题解决后，将关键消息或结论加入收藏，便于后续回看。
\end{usecaselist}

\usecasesubheading{用例二：创建专题聊天室组织集中讨论}

当某个任务需要多人集中协作时，可按以下流程操作：
\begin{usecaselist}
\item 创建聊天室，设置房间名称、技术方向、加入方式和有效期。
\item 邀请相关成员加入房间，围绕同一主题集中讨论并共享文件。
\item 任务结束后保留历史记录，聊天室在到期后转为归档状态。
\end{usecaselist}

\subsection{个人开发者用例分析}

个人开发者的典型操作更偏向先自己验证，再决定是否扩大讨论范围，用例如图~\ref{fig:3-4} 所示。

\begin{figure}[H]
\centering
\thesisfigflow{img/图3_5.png}
\caption{个人开发者用例图}
\label{fig:3-4}
\end{figure}

\usecasesubheading{用例一：向 AI 助手咨询代码问题}

当用户独立开发时，可以通过以下流程获取辅助支持：
\begin{usecaselist}
\item 在 AI 面板中输入代码片段、报错信息或具体问题；遇到陌生术语或局部文本时，也可以直接触发解释。
\item 系统返回解释、修改建议或进一步排查方向，连续追问时保留当前 AI 会话上下文。
\item 将有参考价值的回答保存到收藏页，方便今后再次查看。
\end{usecaselist}

\usecasesubheading{用例二：收藏并整理关键讨论内容}

当用户遇到有长期参考价值的内容时，可按以下流程处理：
\begin{usecaselist}
\item 对目标消息执行收藏操作，并补充必要的备注或标签。
\item 后续在收藏页中按关键词或分类回看对应内容，形成个人资料积累。
\end{usecaselist}

\subsection{普通交流用户用例分析}

普通交流用户的操作相对简单，重点在于快速了解讨论进展和获取相关资料，用例如图~\ref{fig:3-5} 所示。

\begin{figure}[H]
\centering
\thesisfigflow{img/图3_6.png}
\caption{普通交流用户用例图}
\label{fig:3-5}
\end{figure}

\usecasesubheading{用例一：加入聊天室参与日常沟通}

普通交流用户参与房间讨论时，可按以下流程进行：
\begin{usecaselist}
\item 通过公开入口或密码进入目标聊天室，查看公告和历史消息。
\item 根据需要补充问题、确认进度或下载房间中共享的资料文件。
\end{usecaselist}

\usecasesubheading{用例二：检索历史消息并获取阶段性结论}

当用户需要快速了解某个问题的处理结果时，可按以下流程操作：
\begin{usecaselist}
\item 通过历史消息检索或聊天总结功能定位相关讨论内容。
\item 查看总结结果、打开原始消息或下载关联文件，完成信息回读。
\end{usecaselist}

本章从用户、场景、流程、模块和约束几个层面说明了系统需求。Mini-Chat-Bar 的需求重点不只是完成消息收发，还在于让技术讨论中的消息、结论和辅助回答能够被组织、回看和继续利用。后续章节将在此基础上展开系统总体设计与具体实现。
